\documentclass[11pt, a4paper, oneside]{article}
\usepackage[ngerman]{babel}
\usepackage{worksheet}

\begin{document}
	\author{L. Bung}
	\title{Polynomfunktionen: Symmetrie}
	\subject{Mathematik}
	\class{FHR1}
	\maketitle
	
	\partnertask{Untersuchen der Symmetrien}
	
	a) Untersuchen Sie, welche Symmetrien Sie bei den folgenden Funktionen feststellen können.
	Sie können dafür Geogebra verwenden\footnote{\url{https://www.geogebra.org/calculator}}.
	
	\begin{multicols}{3}
		\begin{enumerate}[label=]
			\item $a(x) = \frac{1}{3} x^5 - x^3$
			\item $b(x) = \frac{2}{5} x^4 - 2x^2$
			\item $c(x) = \frac{1}{2} x^4 + \frac{2}{3} x^3$
		\end{enumerate}
	\end{multicols}
	
	\checkered[7cm]
	
	b) Finden Sie ein Muster:
	Welche Symmetrie liegt bei welchen Funktionstermen vor?
	Überprüfen Sie Ihre Hypothese, indem Sie eigene Funktionsterme in Geogebra eingeben.
	
	\checkered[6cm]
	
	\begin{hint}{Symmetrien mathematisch überprüfen}
		Es gibt zwei Arten von Symmetrie:
		\textbf{Achsensymmetrie} und \textbf{Punktsymmetrie}.
	
		\begin{multicols}{2}
			\centering
			\begin{plot}[xmin=-2.5, xmax=2.5, ymin=-1.5, ymax=3.5, restrict y to domain=-5:5]
			\addplot[domain=-2.5:2.5, red, thick, samples=100]{x^4-2*x^2};
			\draw[blue, very thick] (0,-1.5) -- (0,3.25);
			\draw[<->, dashed, blue] (-1.5,1) -- (1.5,1);
			\node[red] at (0,1.5){$f(x)=x^4-2x^2$};
			\end{plot}
			\textbf{Achsensymmetrie}\\
			$f(-x) = f(x)$
			
			\begin{plot}[xmin=-2.5, xmax=2.5, ymin=-2.5, ymax=2.5, restrict y to domain=-5:5]
				\addplot[domain=-2.5:2.5, red, thick, samples=100]{x^5-2*x^3};
				\node[red] at (0,-1.5){$f(x)=x^5-2x^3$};
				\coordinate (S) at (0,0);
				\draw[blue, fill=blue] (S) circle(1.5pt) node[above right, blue]{$S$};
				\draw[<->, dashed, blue] (-1.1,1) -- (1.1,-1);
				\draw[<->, dashed, blue] (-1.5,-1) -- (1.5,1);
			\end{plot}
			\textbf{Punktsymmetrie}\\
			$f(-x)=-f(x)$
		\end{multicols}
	
		Für Polynomfunktionen gilt dabei:
		
		\begin{enumerate}[itemsep=1.5em]
			\item Ein Polynom ist \textit{achsensymmetrisch}, wenn \bigskip
			\item Ein Polynom ist \textit{punktsymmetrisch}, wenn \bigskip
		\end{enumerate}
		
		\phantom{x}
	\end{hint}
	
	
	\singletask{Symmetrien nachrechnen}
	
	a) Geben Sie an, ob die Funktion $f(x) = \frac{1}{3} x^6 - 2x^4 +\frac{1}{2} x$ achsen- oder punktsymmetrisch ist.
	Begründen Sie Ihre Antwort!
	
	\checkered[5.5cm]
	
	b) Rechnen Sie mithilfe der Definition nach, dass die Funktion $g(x) = 2x^5+x^3$ punktsymmetrisch ist.
	
	\checkered[11cm]
	
	c) Können Sie eine Funktion finden, die sowohl punktsymmetrisch als auch achsensymmetrisch ist?
	
	\checkered[8cm]
	
	\let\clearpage\relax
\end{document}